\chapter*{Závěr}
\addcontentsline{toc}{chapter}{Závěr}
\markboth{Závěr}{}

Cílem bakalářské práce bylo analyzovat potenciál vodíku jako nosiče energie a jeho využití v procesu rafinérií, na jeho výrobu v rafinériích \enquote{šedými} i \enquote{zelenými} metodami a především na dekarbonizaci rafinérií, tedy zbavení se závislosti na výrobě vodíku pomocí technologií, které využívají fosilní paliva.

První kapitola je věnována popisu vodíku, jeho fyzikálních vlastností a definicím jeho typů včetně cenového srovnání. 

V druhé kapitole je popsáno jeho využití, zejména v rafinérském průmyslu. Současné využití vodíku a jeho budoucí potenciál se však značně liší. Dnes se vodík převážně používá jako surovina pro výrobu ostatních látek. Očekává se však, že v následujících letech se vodík stane důležitým energetickým nosičem, nejdříve v menších systémech a postupně i v celé energetické infrastruktuře. K dosažení těchto cílů bude nutné výrazně zvýšit produkci vodíku a zároveň splnit evropské cíle pro snižování emisí. \cite{Tkac2017} 

Ve třetí kapitole je popsána výroba vodíku v rafinériích, tedy parním reformingem, parciální oxidací těžkých ropných frakcí a elektrolýzou vody. 

Možnosti dekarbonizace rafinérie jsou popsány ve čtvrté kapitole, věnující se analyzovaným projektům v Evropě. Popisy tamních záměrů ukazují, že vodík může hrát klíčovou roli v dosažení uhlíkové neutrality v rafinérském sektoru. Zatímco některé společnosti vykazují významný pokrok směrem k dekarbonizaci, mnoho projektů je stále ve fázi pilotních studií nebo raných implementací.

Cesta za nejen evropskou nulovou klimatickou stopou je však spojena s významnými výzvami, cíly a náklady. Ty jsou analyzovány v poslední kapitole, která je věnována praktickému výpočtu a finálnímu vyhodnocení dekarbonizace rafinérie. Výpočty ukazují, že při parciální oxidaci těžkých topných olejů vzniká na 1 kilogram vyrobeného vodíku přibližně $12,83$ kilogramů oxidu uhličitého. Nahrazení této technologie výrobou vodíku pomocí elektrolýzy, která by využívala energii z lokální solární elektrárny, by mohlo výrazně snížit emise \ce{CO2}. Navržená solární elektrárna s roční produkcí $164\ 138\ \mathrm{MWh}$ by pokryla až $82\%$ energetických potřeb pro elektrolyzéry, které by nahradily $5\%$ současné výroby vodíku zeleným vodíkem, což by vedlo ke snížení emisí oxidu uhličitého až o $51\ 000$ tun ročně.